\section{Other cultivation circuits}
\label{app:other-cultivation}
\captionsetup{hypcap=false}

We compare other cultivation circuits, ending with a
noiseless code projection and logical readout. Escape and following
stages are omitted for consistency with the main text. We adapt the methods
of the main text to these circuits --- using tensor contraction and compare the resultant LER series with the 
available numerical results.

Recall, $A$ is the probability of acceptance and $B$ is the probability of
acceptance with a logical error, and $\PL=B/A$. Throughout we use this series expansion
\[
 x=\frac{p}{1-p},\qquad S_K(p)=\sum_{k=0}^{K}L_kx^k.
\]
The coefficients below are exact for the specified variants of cultivation. Each protocol's bare distance labels $d=3,5$ must not be confused with the circuit fault distance - $d_\text{fault}$.

\subsection*{1. \texorpdfstring{$\mathbb{RP}^{2}$}{RP2}}

For the actual $T$ converted cultivation circuits from~\cite{RP2_prxq}, we obtain
\begin{equation}
P_{L,\mathbb{RP}^2}^{(3)}(x)=\frac{131}{450}x^{2}+\frac{9874093}{10800}x^{3}+\frac{17172616651}{540000}x^{4}+\frac{27243842363149}{24300000}x^{5}+O(x^{6}).
\end{equation}

\begin{equation}
P_{L,\mathbb{RP}^2}^{(5)}(x)=\frac{10439}{13500}x^{3}+\frac{20300775541}{51840000}x^{4}+\frac{616219992737}{1728000}x^{5}+\frac{428120035103643799}{20736000000}x^{6}+O(x^{7}).
\end{equation}


The leading powers of $x$ establish fault distances $d_{\text{fault}}=2,3$
for $\mathbb{RP}^{2}$ at $d=3,5$, respectively.

\fig{rp2_comparison}{1}{Actual $T$ cultivation series expanded LER and supplied the $S$-gate proxy samples for
$\mathbb{RP}^{2}$~\cite{RP2_prxq}.}

\subsection*{2. Fold-transversal}

For the fold-transversal circuits of~\cite{fold_msc_prxq}, we obtain
\begin{equation}
P_{L,\mathrm{Fold}}^{(3)}(x)=\frac{3909023513}{41674500}x^{3}+\frac{23378177959691}{12602368800}x^{4}+\frac{21954457843901229433}{396974617200000}x^{5}+O(x^{6}).
\end{equation}

\begin{equation}
P_{L,\mathrm{Fold}}^{(5)}(x)=\frac{3493956797}{17781120000}x^{4}+\frac{140969009479148306917}{1411465305600000}x^{5}+\frac{241056568979596536874210373}{32012033131008000000}x^{6}+O(x^{7}).
\end{equation}


The leading powers of $x$ establish fault distances $d_{\text{fault}}=3,4$
for fold-transversal cultivation at $d=3,5$, respectively.

\fig{fold_comparisons}{1}{Published Fold $d=3$ data compared with $B$ per
initial attempt and $\PL=B/A$. The $Y$-proxy curves use independent samples. Black points and
error bars are digitised from the original figures.}

The $Y$-proxy and $T$ data are consistent with $B$ rather than $B/A$. The original no-escape plotting code and counts are unavailable and we cannot confirm the source of the discrepancy. 
It could be possible that the authors were plotting $B$ instead of $B/A$? Or we made a mistake somewhere here.

\subsection*{3. Chan et al.: repaired circuits accompanying arXiv version 1}

For the supplied three-round repaired circuits associated in~\cite{chan2026diagnosingrestoringdegradedfault}, the full-circuit series are
{\small
\begin{equation}
P_{L,\mathrm{Chan},\,v1}^{(3)}(x)=\frac{8514181}{13500}x^{3}+\frac{3637750171}{101250}x^{4}+\frac{9638727601309}{6075000}x^{5}+O(x^{6}).
\end{equation}
\begin{equation}
P_{L,\mathrm{Chan},\,v1}^{(5)}(x)=\frac{53656}{16875}x^{4}+\frac{15796640737639}{24300000}x^{5}+\frac{5280067584675593}{121500000}x^{6}+\frac{22550691826538722573}{10935000000}x^{7}+O(x^{8}).
\end{equation}

}

The leading powers of $x$ establish fault distances $d_{\text{fault}}=3,4$
for Chan et al.'s repaired v1 circuits at $d=3,5$, respectively.

\fig{chan_v1_comparison}{1}{Repaired arXiv v1 circuits: exact-coefficient series
and digitised actual-$T$ sampling. The $d=3$ line retains terms through $x^{10}$ and the
$d=5$ line through $x^7$. These sampling points belong to arXiv v1, not the further
distance-restored four-round circuits in their arXiv v2.}

The four displayed $d=5$ series values lie inside the sampling
uncertainty intervals. For $d=5$, we found an explicit weight-four fault pattern that
passes post-selection and produces a logical error: $X_{10},X_4,X_2$
during the three growth rounds and $Y_3$ during final double checking.
Conditioned on this pattern and no other faults, $A_F=B_F=1/4$.
Separate gate-level ZX calculus contraction reproduce the uncaught logical error.

The following subsection shows results for their arXiv v2 circuits, which add a
fourth stabiliser round.

\subsection*{4. Chan et al.: repaired circuits accompanying arXiv version 2}

The new repaired circuit adds a fourth noisy stabiliser round before final
$d=5$ double checking. The complete series, including injection, both
double-checking stages and all growth noise, is now:
\begin{equation}
P_{L,\mathrm{new}}^{(5)}(x)=\frac{350644373171}{540000}x^5+\frac{15818820811063247}{364500000}x^6+O(x^7).
\label{eq:new-scaling}
\end{equation}

All lower-order coefficients vanish. Thus the circuit fault distance is finally restored to 5.

\begin{center}
\begin{tikzpicture}[font=\scriptsize,line width=.55pt]
\begin{yquant}[register/minimum height=4mm,register/minimum depth=4mm,
 operator/minimum width=7mm,operator/separation=2.3mm,
 every wire/.append style={draw=black,line width=.7pt}]
qubit {$\mathcal D$} q;
[fill=black!4] box {Injection} q;
[fill=paperpurple!8] box {\shortstack{$d=3$\\double checking}} q;
[fill=black!4] box {Expansion} q;
[fill=black!4] box {$R_1$} q;
[fill=black!4] box {$R_2$} q;
[fill=black!4] box {$R_3$} q;
[fill=newround!17,draw=newround,line width=1.1pt] box {$R_4$} q;
[fill=paperpurple!8] box {\shortstack{$d=5$\\double checking}} q;
[fill=black!4] box {\shortstack{Noiseless\\readout}} q;
\end{yquant}
\end{tikzpicture}

\captionof{figure}{The supplied repair. The additional noisy stabiliser
round is green. Boxes group physical gates; the wire represents the data register.}
\end{center}

\fig{chan_four_round_plot_only}{.93}{The new four-round series plotted over the arXiv v1
comparison. Sampling points belong to arXiv v1 only. The inset compares both
circuits truncated at sixth order.}

An independent ZX calculus contraction confirmed rejection of the earlier weight-four fault.
At $p=10^{-3}$, the series through $x^6$ differ by only $0.57\%$, so the
shorter three-round circuit gives nearly the same LER here, although the
fourth round raises the fault distance from $4 \rightarrow 5$, but is somewhat not impactful in the only fruitful figure of merit - LER.

